\section{Reduction to two estimates}

Write \(L=\log x\) and set \(\eta=10^{-6}\).  For a prime \(p\nmid q\), write
\[
        \iota_p(q)=\frac{p-1}{\ord_p(q)}
\]
for the order index of \(q\) modulo \(p\).  The first proposition supplies
primes whose predecessors have one large prime factor and otherwise only
prime factors in a controlled interval.

\begin{proposition}[A prime reservoir in a fixed progression]
\label{prop:reservoir}
Assume Inputs~\ref{input:marked}, \ref{input:block},
and~\ref{input:density}.  Let \(M\ge1\) be odd and squarefree, and let
\(a\) be a residue class modulo \(M\) with
\(\gcd(a(a-1),M)=1\).  There are constants \(c_{M,a}>0\) and \(x_{M,a}>1\)
such that, for every real \(x\ge x_{M,a}\), at least
\(c_{M,a}x/L^2\) primes \(p\) satisfy
\[
 x<p<2x,\qquad p\equiv a\pmod M,\qquad p=4rQ+1,
\]
where \(r\) is a positive integer, \(Q>x^{0.9}\) is prime, and
\[
 \exp(L^{0.1})<\ell<\exp(L^{0.3})
 \qquad\text{for every prime }\ell\mid r.
\]
The case \(M=1\) is included.
\end{proposition}

The second proposition bounds the exceptional primes at which a controlled
prime factor can divide the index of a prescribed base.

\begin{proposition}[A uniform bound for order obstructions]
\label{prop:splitting}
Assume Input~\ref{input:hecke}.  For each fixed positive prime \(q\), for all
sufficiently large real \(x\), and uniformly over primes \(\ell\) with
\(\exp(L^{0.1})<\ell<\exp(L^{0.3})\), one has
\[
 \#\left\{\,p\text{ prime}:
 \begin{array}{l}
 x<p<2x,\quad p\nmid q\ell,\quad p\equiv1\pmod\ell,\\
 q^{(p-1)/\ell}\equiv1\pmod p
 \end{array}\right\}
 \ll_q \frac{x}{\ell(\ell-1)L}+x^{1-\eta}.
\]
\end{proposition}

\begin{proof}[Proof of Theorem~\ref{thm:main}]
Let \(M\) be the product of the odd members of
\(\{q_1,\ldots,q_k\}\), with \(M=1\) if this set has no odd member.
For each odd \(q_i\), choose a nonzero quadratic nonresidue modulo \(q_i\).
The Chinese remainder theorem supplies a class \(a\bmod M\) having all
these prescribed residues; when \(M=1\), take its unique residue class.
Since \(1\) is a square, \(\gcd(a(a-1),M)=1\).

Apply Proposition~\ref{prop:reservoir}, and call the resulting set of
primes \(\mathcal R(x)\).  For large \(x\), all its primes exceed the
\(q_i\), and both \(r\) and \(Q\) in \(p=4rQ+1\) are odd.
Thus \(p\equiv5\pmod8\), in particular \(p\equiv1\pmod4\).
Quadratic reciprocity and the supplementary law for \(2\) give
\[
 \left(\frac{q_i}{p}\right)
 =
 \begin{cases}
 \displaystyle\left(\frac{p}{q_i}\right)
   =\left(\frac{a}{q_i}\right)=-1,&q_i>2,\\[6pt]
 \displaystyle(-1)^{(p^2-1)/8}=-1,&q_i=2.
 \end{cases}
\]
For a fixed one of the bases \(q=q_i\), the index \(\iota_p(q)\) is odd:
if it were even, then
\(\ord_p(q)\mid(p-1)/2\), contrary to Euler's criterion.
If \(\iota_p(q)>1\), it has an odd prime divisor \(\lambda\).  Since
\(\iota_p(q)\mid4rQ\), either \(\lambda=Q\) or \(\lambda\mid r\), and
in either case
\begin{equation}\label{eq:index-obstruction}
        q^{(p-1)/\lambda}\equiv1\pmod p.
\end{equation}

If \(Q\mid\iota_p(q)\), the integer \(j=(p-1)/Q\) is less than
\(2x^{0.1}\), and \(p\mid q^j-1\).  The product of \(q^j-1\) over
\(1\le j\le2x^{0.1}\) has logarithm at most
\[
        (\log q)\sum_{j\le2x^{0.1}}j\ll_q x^{0.2}.
\]
Each of its prime divisors in \((x,2x)\) contributes more than \(L\) to
this logarithm.  Consequently the number of these exceptional primes is
\begin{equation}\label{eq:large-factor-exceptions}
        O_q\!\left(\frac{x^{0.2}}{L}\right)
        =o\!\left(\frac{x}{L^2}\right).
\end{equation}

In the other case, write \(\ell=\lambda\mid r\).  The controlled interval
in Proposition~\ref{prop:reservoir} applies to \(\ell\).  Since
\(\ell\mid p-1\), one has \(p\equiv1\pmod\ell\), while
Equation~\eqref{eq:index-obstruction} supplies the other congruence in
Proposition~\ref{prop:splitting}.  Sum that proposition
over all such primes \(\ell\).  The sum of \(1/(\ell(\ell-1))\) is at most
the telescoping sum over all integers exceeding \(\exp(L^{0.1})\), and
there are at most \(\exp(L^{0.3})\) possible \(\ell\).  The number of
exceptions of this second kind is therefore
\begin{equation}\label{eq:controlled-factor-exceptions}
 O_q\!\left(
       \frac{x}{L}\exp(-L^{0.1})
       +\exp(L^{0.3})x^{1-\eta}
       \right)
 =o\!\left(\frac{x}{L^2}\right).
\end{equation}
Here the two ratios to \(x/L^2\) are bounded by constant multiples of
\(L\exp(-L^{0.1})\) and
\(L^2\exp(L^{0.3}-\eta L)\), respectively, and both tend to zero.

There are only finitely many bases.  The union over them of the exceptions
in Equations~\eqref{eq:large-factor-exceptions} and
\eqref{eq:controlled-factor-exceptions} is still \(o(x/L^2)\).  Since
\(\#\mathcal R(x)\ge c_{M,a}x/L^2\), for every sufficiently large real
\(x\) at least \(c_{M,a}x/(2L^2)\) primes in \(\mathcal R(x)\) have
\(\iota_p(q_i)=1\) for every \(i\).  These primes give the theorem.
\end{proof}
